\section{پاسخ سوال ۸}
\begin{stepbox}
\field{محاسبه ماتریس ساختار در روش هریس (Harris Corner)}{
پیکسل مرکزی در بلاک ۵×۵ داده شده، در سطر سوم و ستون سوم با مقدار $110$ قرار دارد. برای محاسبه ماتریس $M$، به گرادیان‌های مکانی $I_x$ و $I_y$ در یک پنجره ۳×۳ حول این پیکسل نیاز داریم. 

\textbf{محاسبه گرادیان‌ها با تقریب تفاوت مرکزی:}
گرادیان افقی $I_x(x,y) = I(x, y+1) - I(x, y-1)$ و گرادیان عمودی $I_y(x,y) = I(x+1, y) - I(x-1, y)$.
با اعمال این رابطه برای تمام ۹ پیکسل موجود در همسایگی ۳×۳، مقادیر زیر حاصل می‌شود:

\begin{codebox}[خروجی مقادیر مشتقات در پایتون]
Ix = [[  7,   3,  -7], 
      [ 52,  49, -10], 
      [ 54,  -4,  -4]]

Iy = [[ 19,  67,  66], 
      [ 66,  57,  59], 
      [ 61,   7,   7]]
\end{codebox}

\textbf{محاسبه درایه‌های ماتریس $M$:}
ماتریس ساختار (Structure Tensor) برابر است با مجموع ضرب‌های گرادیان در پنجره (بدون وزن‌دهی گاوسی برای سادگی آزمون):
\[ M = \sum \begin{bmatrix} I_x^2 & I_x I_y \\ I_x I_y & I_y^2 \end{bmatrix} \]
با به توان دو رساندن و ضرب المان به المان مقادیر فوق و جمع آن‌ها در کل پنجره ۳×۳:
\[ \sum I_x^2 = (7)^2 + (3)^2 + (-7)^2 + (52)^2 + (49)^2 + (-10)^2 + (54)^2 + (-4)^2 + (-4)^2 = \textbf{8260} \]
\[ \sum I_y^2 = (19)^2 + (67)^2 + (66)^2 + (66)^2 + (57)^2 + (59)^2 + (61)^2 + (7)^2 + (7)^2 = \textbf{24111} \]
\[ \sum I_x I_y = (7\times19) + (3\times67) + \dots + (-4\times7) = \textbf{8745} \]

در نهایت ماتریس $M$ به صورت زیر به دست می‌آید:
\[ M = \begin{bmatrix} 8260 & 8745 \\ 8745 & 24111 \end{bmatrix} \]

\begin{tipbox}[تأیید پیاده‌سازی با ابزارهای تخصصی]
در توسعه الگوریتم‌های بینایی ماشین، می‌توان این عملیات را در کتابخانه‌های تخصصی نظیر \lr{OpenCV} (تابع \texttt{cornerHarris}) یا مدل‌های شبکه‌های عصبی پیچشی (با استفاده از \lr{PyTorch} که در آن فیلترهای مشتق‌گیر به عنوان وزن‌های کرنل تعریف می‌شوند) با سرعت بسیار بالایی در قالب عملیات تانسوری پیاده‌سازی کرد.
\end{tipbox}
}
\end{stepbox}
